%% ============================================================================
%% coverpage.tex -- preliminary pages א׳–ד׳ (booklet Table 1.1, chapter 2)
%%   א׳ cover (no visible page number) | ב׳ executive summary + keywords
%%   ג׳ table of contents (depth 3)    | ד׳ lists of figures, tables, abbreviations
%% The executive summary is written LAST (booklet 2.2): 250-350 words, one block.
%% ============================================================================
\makecoverpage

%% ------------------------------------------------------------------ page ב׳
\chapter*{תקציר מנהלים}
\markboth{תקציר מנהלים}{}
\begingroup\setstretch{1.2}
ענף הביוטי והאופנה המקוון נשען על החלטות של התאמה אישית: גוון מייקאפ, צבע שיער ומידת בגד. בפועל, תמונת סלפי מצולמת בתאורה לא מבוקרת, ומערכות ההמלצה הקיימות מקבלות נתון צבע מעוות ונתון מידות חסר. סימולציה ספקטרלית של הצוות על \num{4,000} סלפי סינתטיים, בארבע משפחות תאורה, הראתה שגיאת צבע ממוצעת של \num{$\Delta E_{00} = 6.58$} בעור ללא תיקון, ושאיזון לבן גלובלי מגדיל אותה ל-\num{16.03}, הרבה מעל סף ההבחנה של כ-\num{2.0}. במקביל, הדמיה גנרטיבית בשירות מסחרי עולה \num{\$0.045--0.090} לתמונה, בלי מנגנון שמממן אותה.
הדוח מציע את \en{AURAM} (\en{Aura} ו-\en{RAM}), מערכת רב-מודאלית שמכיילת את צבעי המשתמשת בשני אומדנים בלתי תלויים: עוגן ביומטרי בעין, לובן העין שנמדד ברישום וצבע העיניים המדווח כהתפלגות קודמת בייסיאנית, והפרש בין תמונה עם מסך לבן לתמונה בלעדיו, שמבודד את ההחזרה מהאור. אל הכיול מצטרפים אומדן היקפי גוף במודל \en{SMPL-X} בעזרת הגובה המדווח, מנוע אחזור רב-מודאלי (\en{RAG}) שמחבר קטלוג מוצרים לטרנדים ברשתות החברתיות, ומודל הדמיה מזוקק שממומן במנוע החלפת ערך בשלוש שכבות.
יעדי התכן המרכזיים הם חציון \num{$\Delta E_{00} \le 1.8$}, שגיאת היקף \num{$\text{MAE} \le 1.2$ cm}, זמן המלצה \num{$\le 1.8$ s}, הדמיה \num{$\le 4.5$ s}, והכנסה של פי \num{2.4} לפחות מעלות החישוב. שלוש ארכיטקטורות נבחנו במטריצת החלטה משוקללת שמשקליה נקבעו מראש. הארכיטקטורה ההיברידית של חישוב מפוצל קיבלה \num{4.30}, מול \num{3.40} לעיבוד בקצה ו-\num{2.85} לענן מונוליטי, ונשארה מובילה בבדיקות הרגישות למשקלים ולציונים. בבדיקת היתכנות סימולטיבית הוריד הכיול המוצע את האחוזון ה-\num{95} ל-\num{4.25}, מתחת ליעד, אך החציון, \num{1.86}, חורג מעט מהיעד של \num{1.8}. מודל בייסיאני היררכי נותן לחציון רווח בר-סמך של \num{1.75--2.00} וסיכוי של \num{12\%} בלבד לעמוד ביעד בתאורות שנבדקו, ומראה שמשפחות תאורה נוספות שוות לניסוי יותר ממשתתפות נוספות. שתי השערות ייבחנו במודל מעורב בניסוי של \num{120} משתתפות. בהנחות צוות שטרם אומתו, הנתח הניתן להשגה בשנה השלישית נאמד בכ-\num{\$45.0M}, ו-\en{LTV/CAC} בחציון של \num{6.1}, בטווח \num{2.2--18.3}, עם סיכוי של \num{78\%} לעבור את סף היעד; ניתוח ערך המידע מצביע על ניסוי ההמרה אצל השותף כמדידה הבאה. כל התוצאות מופקות מקוד שזמין במאגר הפרויקט וניתנות לשחזור.
המימוש מתוכנן ל-\num{14} שבועות בחמש חבילות עבודה, עם נתיב קריטי, שתי אבני דרך בשבועות \num{8} ו-\num{12}, ושמונה סיכונים מכומתים עם תוכנית הפחתה.
\par\endgroup
\vspace{8pt}
\begin{keymessagebox}[מילות מפתח]
דוח מכין; כיול ביומטרי רב-מודאלי; קביעות צבע במרחב \en{CIELAB}; רשת גוף תלת-ממדית; היווצרות תמונה ספקטרלית; מודל בייסיאני היררכי; ערך המידע; הדמיה וירטואלית גנרטיבית; כלכלת החלפת ערך בשכבות.
\end{keymessagebox}
\vfill
\begin{mirrorbox}[מראת הכיס: חתימה ולא קישוט]
מראת כיס היא חפץ שכמעט לא השתנה במאה שנים, ובכל זאת כל דור עיצב אותה מחדש: כסף חרוט, אמייל, בקליט, זהב ורוד. הצללית והגימור מתחלפים, והיא מזוהה תמיד במבט אחד. זו גם ההבטחה של \en{AURAM}: הטרנדים מתחלפים בתוך שבועות, והכיול האישי נשאר. מראת כיס אינה מייפה ואינה שופטת. שולפים אותה לרגע, היא עונה על שאלה אחת, איך אני נראית עכשיו ובאור הזה, ונסגרת; כך בדיוק אמורה לפעול שכבה שקופה (סעיף~\ref{sec:human}). ומי שנושאת מראה בתיק כבר החליטה שהמראה שלה הוא בחירה ולא מקרה. לכן הסמל חתום מתחת לכל מספר עמוד בדוח: קטן, שקט ותמיד שם.
\end{mirrorbox}

%% ------------------------------------------------------------------ page ג׳
\compacttoc

%% ------------------------------------------------------------------ page ד׳
\compactlists{%
\begin{abbrevtable}
\abbr{API}{ממשק תכנות יישומים}
\abbr{AR}{מציאות רבודה}
\abbr{ARPU}{הכנסה ממוצעת למשתמשת}
\abbr{CAGR}{שיעור צמיחה שנתי מורכב}
\abbr{CIELAB}{מרחב צבע אחיד-תפיסתית}
\abbr{CLIP}{מודל שיכון משותף לתמונה וטקסט}
\abbr{$\Delta E_{00}$}{מדד הפרש צבע \en{CIEDE2000}}
\abbr{eCPM}{הכנסה לאלף צפיות}
\abbr{EVPI}{ערך צפוי של מידע מושלם}
\abbr{FAISS}{ספריית חיפוש וקטורי}
\abbr{GDPR}{תקנת הגנת המידע האירופית}
\abbr{GPU}{מעבד גרפי}
\abbr{IoU}{חפיפה יחסית}
\abbr{LLM}{מודל שפה גדול}
\abbr{MAE}{שגיאה מוחלטת ממוצעת}
\abbr{PRD}{מסמך דרישות מוצר}
\abbr{RACI}{מטריצת אחריות: מבצע, אחראי, נועץ, מיודע}
\abbr{RAG}{יצירה מועשרת-אחזור}
\abbr{SAM}{השוק הזמין לשירות}
\abbr{SKU}{יחידת מלאי קטלוגית}
\abbr{SMART}{יעד ספציפי, מדיד, בר-השגה, רלוונטי ותחום בזמן}
\abbr{SMPL-X}{מודל גוף פרמטרי תלת-ממדי עם ידיים ופנים}
\abbr{SOM}{הנתח הניתן להשגה}
\abbr{TAM}{השוק הכולל}
\abbr{VTON}{הדמיה וירטואלית של לבוש}
\abbr{WBS}{מבנה פירוק עבודה}
\end{abbrevtable}}
